% AAAI 2026 Submission — Anonymous
\documentclass[letterpaper]{article}
\usepackage{aaai26}
\usepackage{times}
\usepackage{helvet}
\usepackage{courier}
\usepackage[T1]{fontenc}
\usepackage{url}
\usepackage{graphicx}
\usepackage{latexsym}
\usepackage{amsmath}
\usepackage{amssymb}
\usepackage{booktabs}
\usepackage{multirow}
\usepackage{makecell}
\usepackage{xcolor}
\usepackage{pifont}
\usepackage{microtype}
\usepackage{inconsolata}

\pdfinfo{
  /Title (Negation Blindness in Large Language Models: Diagnosis and Multi-Grained Training Mitigation)
  /Author (Anonymous)
}

\setcounter{secnumdepth}{2}
\setlength\titlebox{2.5in}

\newcommand{\cmark}{\ding{51}}
\newcommand{\xmark}{\ding{55}}
\newcommand{\suppress}{[\textsc{Suppress}]}
\newcommand{\preserve}{[\textsc{Preserve}]}
\newcommand{\select}{[\textsc{Select}]}
\newcommand{\mgnm}{\textsc{mgnm}}
\newcommand{\spos}{s(\cdot)}

\title{Negation Blindness in Large Language Models:\\
Diagnosis and Multi-Grained Training Mitigation}

\author{Anonymous}

\begin{document}
\maketitle

%% ─────────────────────────────────────────────────────────────────────────────
\begin{abstract}
Large language models (LLMs) exhibit \emph{negation blindness}: when
a prompt explicitly forbids a concept, models frequently output that
very concept anyway.
We make three contributions.
\textbf{(1)} We release \textbf{E4}, a 937-record benchmark covering three
negation phenomena---target suppression, scope preservation, and
negated ranking---with 99.6\% entity holdout.
\textbf{(2)} We present \textbf{MGNM} (Multi-Grained Negation Modeling), a
LoRA fine-tuning framework that uses caller-specified behavior tokens with
five complementary loss terms: suppression, ranking, preservation, positive
retention, and capability retention.
\textbf{(3)} On a 397-record label-audited E4 v2 test subset, MGNM achieves
64.2\% \textsc{NegationFlipAcc} (\textsc{NegFlipAcc}) and 92.9\% \textsc{ScopeCtrl} on Qwen2.5-7B,
outperforming all baselines on both metrics simultaneously
(vs.\ NC-SFT: $+13.4$pp NegFlipAcc and $+15.3$pp ScopeCtrl).
On BoolQ, Qwen MGNM reaches 87.2\% after PMI calibration, while Llama
exposes a No-prior calibration shift and training tradeoff: its raw score of
69.0\% recovers to 81.4\% under standard PMI.
The same trend holds on Llama-3.1-8B, with Mistral-7B providing partial
cross-architecture validation.
These peak E4 results are conditional on the correct behavior token;
structure-assisted behavior routing is evaluated separately as a
deployment-calibration setting under known task structure, not as a solved
pure-text automatic routing problem.
Mechanistic probes on 200 E4 candidate-ranking instances suggest that,
under this setting, negation failures are not primarily caused by lack
of attention to negation tokens, but by a failure to translate the
negation signal into candidate-level ranking preferences.
MGNM partially repairs this internal signal and provides a discrete interface
for selecting the intended negation behavior.
\end{abstract}

%% ─────────────────────────────────────────────────────────────────────────────
\section{Introduction}
\label{sec:intro}

Consider prompting a generative language model with:
\begin{quote}
  \emph{``Name a noun that does \textbf{not} belong to the category of birds.''}
\end{quote}
A well-calibrated model should output a non-bird entity.
In practice, contemporary LLMs frequently output
\emph{sparrow} or \emph{eagle} --- the very concepts the negation was intended
to exclude.
We term this failure mode \emph{negation blindness}.

Negation blindness is not merely a curiosity.
In retrieval, summarization, and instruction following, negation is a
primary tool for expressing constraints: ``do not include X,''
``avoid Y,'' ``explain what is \emph{not} a valid approach.''
A model that ignores negation produces outputs that violate user intent in
ways that may be difficult to detect.

Prior work has documented that language models struggle with negation in
probe studies~\cite{kassner2020negated,ettinger2020bert},
natural language inference~\cite{hossain2020analysis}, and
reading comprehension~\cite{ravichander2022condaqa}.
However, most studies either (a) treat negation handling as a static
capability measurable via accuracy on a fixed benchmark, without
distinguishing the \emph{type} of failure, or (b) attempt prompt-level
interventions that do not target the root cause~\cite{jang2023negation}.

We make the following contributions:

\begin{itemize}
\item \textbf{E4 Benchmark} (\S\ref{sec:benchmark}): 937 records covering
  three negation phenomena at entity-level granularity, with
  strict entity and template holdout to prevent surface-form leakage.
\item \textbf{MGNM} (\S\ref{sec:method}): a LoRA training framework with
  caller-specified behavior tokens and five complementary losses that jointly
  train suppression, scope-sensitivity, and ranking in one pass.
\item \textbf{Comprehensive evaluation} (\S\ref{sec:experiments}): eight
  baselines, three architectures, two out-of-domain tasks (WikiFact, BoolQ),
  McNemar significance tests, and mechanistic probing.
\item \textbf{Mechanistic diagnosis} (\S\ref{sec:mechanism}): logit-lens,
  attention, and activation-patching experiments on 200 E4 instances suggest that,
  under the candidate-ranking setting, negation failure is not primarily an
  attention deficiency but a failure to translate the negation signal into
  candidate-level ranking preferences.
\end{itemize}

%% ─────────────────────────────────────────────────────────────────────────────
\section{Related Work}
\label{sec:related}

\paragraph{Negation in NLP.}
Negation has long been a challenge in NLP, with task-specific work on
negation detection and scope resolution~\cite{hossain2020analysis}.
In the era of pre-trained models, \citet{kassner2020negated} showed that BERT
treats negated probes almost identically to their affirmative counterparts.
\citet{hosseini2021understanding} proposed negative training examples to
improve negation consistency in masked LMs.
Recent work~\cite{nos2024negation,truong2023unsee} extends these findings to
instruction-tuned LLMs, showing that explicit chain-of-thought improves but
does not solve negation handling.

\paragraph{Fine-tuning for instruction following.}
Supervised fine-tuning~\cite{ouyang2022training} and DPO~\cite{rafailov2023direct}
are the dominant paradigms for aligning LLMs.
Neither is designed for negation-specific constraints.
Vanilla SFT on negation examples leads to over-suppression
(our audited evaluation confirms ScopeCtrl drops to 54.1\%), and DPO's
pairwise objective does not distinguish \emph{suppress} from \emph{preserve} contexts.

\paragraph{Inference-time methods.}
Contrastive Decoding~\cite{li2023contrastive} amplifies the expert--amateur
logit gap to improve output quality.
We show that it exacerbates negation blindness (ScopeCtrl $-16.4$pp)
because the logit gap reflects semantic plausibility, not negation scope.
Chain-of-thought~\cite{wang2022self} and persona prompting fail similarly:
surface-level reasoning does not alter intermediate representations.

\paragraph{Constrained text generation.}
Lexical constraint methods such as Constrained Beam Search~\cite{anderson2017guided}
and NeuroLogic Decoding~\cite{lu2021neurologic} allow specifying inclusion or
exclusion constraints over output tokens.
Our task shares the spirit of exclusion constraints (``do not output $c^*$''),
but operates on \emph{semantic} categories rather than surface tokens, and requires
the model to maintain correct scope discrimination (Preserve behavior) simultaneously.

\paragraph{Instruction-following and constraint satisfaction.}
Recent benchmarks such as IFEval~\cite{zhou2023ifeval} and
FollowBench~\cite{jiang2023followbench} evaluate whether instruction-tuned
LLMs satisfy verifiable or multi-level constraints.
These works measure broad constraint satisfaction as a behavioral property
across diverse instruction types.
E4 is complementary: it isolates negation-specific constraint failures
and separates suppress, preserve, and select behaviors that broad benchmarks
conflate, enabling fine-grained diagnosis of \emph{why} models fail under
negated constraints rather than \emph{whether} they do.

\paragraph{Calibration and prior shift.}
\citet{zhao2021calibrate} show that few-shot prompts introduce spurious label
priors in language models, which can be corrected by subtracting null-prompt scores.
\citet{holtzman2021surface} identify similar prior-shift artifacts as
\emph{surface form competition} and show PMI scoring mitigates them.
We observe that MGNM's $\mathcal{L}_\text{sup}$ induces a null-context prior shift
($\Delta\,{\approx}\,{+1.0}$) that contributes to raw BoolQ degradation; PMI
calibration recovers much of this loss, while retention checkpoints show a
remaining training tradeoff with E4 behavior control.

\paragraph{Representation editing and behavioral steering.}
Li et al.~\cite{li2023inference} show that a small set of attention heads
can be steered at inference time (\emph{inference-time intervention}) to improve
truthfulness.  Turner et al.~\cite{turner2023activation} demonstrate that
adding a fixed steering vector to residual activations consistently shifts model
behavior.  These methods alter continuous representations and cannot distinguish
negation behaviors across record types.  Our behavior-condition tokens
(\textsc{[Suppress]}/\textsc{[Preserve]}) function as a discrete steering
interface: when the intended behavior is supplied, the same negated prompt can
be routed to type-appropriate behavior without modifying base representations.

\paragraph{Mechanistic interpretability.}
Logit-lens probing~\cite{nostalgebraist2020logitlens} and activation
patching~\cite{wang2022interpretability,meng2022locating} have revealed
how factual knowledge is stored in transformer layers.
We apply these methods to diagnose \emph{where} in the processing pipeline
negation information is lost, finding that the failure is not attentional
but representational: negation tokens are attended to but their signal is
not translated into candidate-level ranking preferences.

%% ─────────────────────────────────────────────────────────────────────────────
\section{The E4 Benchmark}
\label{sec:benchmark}

\subsection{Task Formulation}
Given a \emph{negated} prompt $p^-$ and a candidate pool
$\mathcal{C} = \{c_1, \ldots, c_K\}$, a model must rank candidates by
assigning log-probabilities $s(p^-, c_i)$.
Correct behavior depends on the negation \emph{type}:

\begin{description}
\item[\textbf{Suppress}:] The target $c^* \in \mathcal{C}$ is the affirmative
  answer and should rank \emph{last} after negation.
  \emph{``Name a noun that does \textbf{not} belong to birds''} $\to$
  \emph{sparrow} must be beaten by all non-bird alternatives.
\item[\textbf{Preserve}:] The affirmative answer remains correct despite
  surface-level negation in the prompt.
  \emph{``You must \textbf{not} neglect the workflow rules''} $\to$
  the correct action (\emph{obtain approval first}) is still best.
\item[\textbf{Select}:] A negated alternative exists and should rank highest.
  \emph{``Steps that should \textbf{not} be used''} $\to$
  an \emph{incorrect} workflow step should outrank the correct one.
\end{description}

\subsection{Construction and Statistics}
E4 records are generated from 47 templates spanning 6 domains (natural
science, office workflow, geography, social norms, cooking, sports).
Each template instantiates 20 entity families; entities within a family
share semantic structure but differ in surface form.
We enforce strict holdout: test entities share no family with training entities
(holdout rate 99.6\%) and all templates are split at family level.

The full generated dataset contains 937 records (412 train, 525 test v2).
The canonical training split contains 83 suppress, 123 preserve
(out-of-scope), 74 preserve (double-negation), and 132 select instances.
Table~\ref{tab:dataset} summarizes the split statistics.
After a label audit, we remove direct negated-selection prompts that had been
assigned to non-select behaviors (e.g., ``Which planet does not have rings?''
with a ringed-planet positive answer), plus preserve rows whose negated prompt
directly asks for the non-positive answer. Unless otherwise stated, E4 v2
candidate-ranking results are reported on this 397-record audited subset
(97 suppress, 98 preserve, 202 select).
The strongest Qwen adapter reported below is trained on a 471-record augmented
input that adds 59 exclusive-choice select records to the canonical training
split and is evaluated on the audited test subset; the strict audited training
split contains 357 records and is used for audit/phase-2 diagnostics, not as
the source of the strongest adapter.

\begin{table}[t]
\centering
\small
\caption{E4 dataset statistics.}
\label{tab:dataset}
\begin{tabular}{lrrr}
\toprule
Type & Train & Test & Total \\
\midrule
Suppress (in-scope) & \phantom{0}83 & 129 & 212 \\
Preserve (out-of-scope) & 123 & 128 & 251 \\
Preserve (double-neg.) & \phantom{0}74 & \phantom{0}66 & 140 \\
Select (gold-neg) & 132 & 202 & 334 \\
\midrule
\textbf{Total} & \textbf{412} & \textbf{525} & \textbf{937} \\
\bottomrule
\end{tabular}
\end{table}

%% ─────────────────────────────────────────────────────────────────────────────
\section{Method: MGNM}
\label{sec:method}

\subsection{Behavior-Conditioned Prompts}

For each training record with negated prompt $p^-$, we prepend a
\emph{behavior token} that signals the required response type:

\begin{equation}
\tilde{p}^- = \begin{cases}
  \texttt{[SUPPRESS]}\;p^- & \text{if suppress} \\
  \texttt{[PRESERVE]}\;p^- & \text{if preserve} \\
  p^- & \text{if select}
\end{cases}
\end{equation}

The same token is prepended at evaluation time, allowing the model to
condition on explicit behavioral intent.
This is analogous to control codes~\cite{ouyang2022training} but operates
at the negation-type level rather than the style level.

\subsection{Loss Functions}

Let $s(p, c) = \sum_t \log p_\theta(c_t \mid p, c_{<t})$ be the
sequence log-probability of completion $c$ given prefix $p$.
For each record we compute a weighted sum of five losses.

\paragraph{Positive retention ($\mathcal{L}_\text{pos}$).}
Ensures the model preserves affirmative-condition accuracy on the
positive prompt $p^+$:
\begin{equation}
\mathcal{L}_\text{pos} = \operatorname{ReLU}\!\left(
  \gamma_\text{pos} - \max_{c \in \mathcal{G}^+} s(p^+, c)
  + \max_{c \in \mathcal{C}^-} s(p^+, c)
\right)
\end{equation}
where $\mathcal{G}^+$ is the gold affirmative set and
$\mathcal{C}^-$ the competitor pool.

\paragraph{Suppression ($\mathcal{L}_\text{sup}$).}
For \textsc{suppress} records, trains the model to rank all pool
candidates above the forbidden target under the negated prompt:
\begin{equation}
\mathcal{L}_\text{sup} = \operatorname{ReLU}\!\left(
  \gamma_\text{sup} + \max_{c \in \mathcal{G}^+} s(\tilde{p}^-, c)
  - \operatorname{mean}_{c \in \mathcal{P}} s(\tilde{p}^-, c)
\right)
\end{equation}
where $\mathcal{P}$ is the candidate pool (non-target completions).
Note that $\mathcal{L}_\text{sup}$ uses a mean-pool margin; NegFlipAcc strictly
requires \emph{every} pool candidate to beat the target, so a pairwise version
$\sum_{c \in \mathcal{P}} \operatorname{ReLU}(\gamma + s(\tilde{p}^-, c^*) - s(\tilde{p}^-, c))$
is theoretically more aligned.
Preliminary experiments found the mean-pool margin empirically sufficient;
pairwise and softmin variants yielded no statistically significant gain on
the E4 dev set.

\paragraph{Ranking ($\mathcal{L}_\text{rank}$, $\mathcal{L}_\text{dist}$).}
For \textsc{select} records, first trains the negated gold above the
affirmative gold, then pushes it above all distractors:
\begin{align}
\mathcal{L}_\text{rank} &= \operatorname{ReLU}\!\left(
  \gamma - \max_{c \in \mathcal{G}^-} s(p^-, c)
  + \max_{c \in \mathcal{G}^+} s(p^-, c)
\right) \\
\mathcal{L}_\text{dist} &= \operatorname{ReLU}\!\left(
  \gamma - \max_{c \in \mathcal{G}^-} s(p^-, c)
  + \max_{c \in \mathcal{D}} s(p^-, c)
\right)
\end{align}
where $\mathcal{G}^-$ is the gold negated set and $\mathcal{D}$ the distractor pool.

\paragraph{Preservation ($\mathcal{L}_\text{pre}$).}
For \textsc{preserve} records, the affirmative answer must remain top-ranked
even under the negated prompt:
\begin{equation}
\mathcal{L}_\text{pre} = \operatorname{ReLU}\!\left(
  \gamma_\text{pre} - \max_{c \in \mathcal{G}^+} s(\tilde{p}^-, c)
  + \max_{c \in \mathcal{C}^-} s(\tilde{p}^-, c)
\right)
\end{equation}

\paragraph{Capability retention ($\mathcal{L}_\text{ret}$).}
KL divergence between the fine-tuned and base model on a set of
general-domain retention texts $\mathcal{T}$:
\begin{equation}
\mathcal{L}_\text{ret} = \mathbb{E}_{t \in \mathcal{T}}\!\left[
  \mathrm{KL}\!\left(p_\theta(\cdot \mid t) \;\|\; p_{\theta_0}(\cdot \mid t)\right)
\right]
\end{equation}

\paragraph{Total loss.}
\begin{equation}
\mathcal{L} = \lambda_\text{pos}\mathcal{L}_\text{pos}
+ \lambda_\text{sup}\mathcal{L}_\text{sup}
+ \lambda_\text{rank}\mathcal{L}_\text{rank}
+ \lambda_\text{dist}\mathcal{L}_\text{dist}
+ \lambda_\text{pre}\mathcal{L}_\text{pre}
+ \lambda_\text{ret}\mathcal{L}_\text{ret}
\end{equation}

\subsection{Training Setup}

We fine-tune LoRA adapters (rank 16, $\alpha=32$, dropout=0.05,
applied to all seven projection matrices) for 3 epochs using the Adam
optimizer ($\mathrm{lr}=2\times10^{-4}$ for Qwen/Llama,
$5\times10^{-5}$ for Mistral; gradient accumulation $=8$).
Loss weights: $\lambda_\text{pos}=1.0$, $\lambda_\text{sup}=1.0$,
$\lambda_\text{rank}=1.5$, $\lambda_\text{dist}=3.0$,
$\lambda_\text{pre}=1.5$, $\lambda_\text{ret}=0.05$.
Preserve-type records are oversampled $3\times$ to compensate
for class imbalance.
Training requires a single 80GB A100; wall-clock time is 15--40 minutes.

%% ─────────────────────────────────────────────────────────────────────────────
\section{Experiments}
\label{sec:experiments}

\subsection{Metrics}
\label{sec:metrics}

\begin{itemize}
\item \textbf{NegationFlipAcc (NegFlipAcc)}: fraction of flip-required records
  (\texttt{suppress\_target} and \texttt{select\_gold\_neg}) where the model is
  correct on both the positive prompt and the negated prompt. On the audited
  subset this metric has $n{=}299$ records (97 suppress + 202 select).
\item \textbf{ScopeCtrl}: fraction of preserve records where the model
  avoids over-negation (affirmative answer remains top-ranked; $n{=}98$).
\item \textbf{Hard NegRank}: fraction of \texttt{select\_gold\_neg} records
  where the original single gold negative answer is top-ranked ($n{=}202$).
\item \textbf{Multi-answer NegRank}: an audited SELECT metric that also accepts
  manually verified semantically valid negative alternatives when a prompt
  admits more than one correct negative answer.
\item \textbf{OverNeg}: complement of ScopeCtrl (lower is better).
\item \textbf{MMLU $\Delta$}: accuracy change on a 570-question MMLU 5-shot
  diagnostic subset relative to the base model, used as a lightweight
  capability-retention check.
\end{itemize}

\subsection{Baselines}

\paragraph{Inference-time.}
\emph{Prompt+Warning}: system message instructing the model to strictly respect negation.
\emph{Prompt+Persona}: persona prompt framing the model as a negation-aware expert.
\emph{Prompt+CoT}: five-step chain-of-thought scaffold.
\emph{Contrastive Decoding} (CD)~\cite{li2023contrastive}: Qwen2.5-7B (expert)
minus $0.5\times$ Qwen2.5-0.5B (amateur) log-probabilities.

\paragraph{Fine-tuning.}
\emph{Vanilla SFT}: standard supervised fine-tuning on E4 training records
with cross-entropy on the gold completion.
\emph{DPO}~\cite{rafailov2023direct}: direct preference optimization with
negated completions as rejected and gold completions as chosen.
\emph{NC-SFT}: negation-conditioned SFT using only cross-entropy on negated
examples (no behavior tokens, no preservation loss).
Fine-tuned main-table adapters follow the corresponding training configuration;
the strongest Qwen MGNM adapter uses the 471-record augmented training input
described in \S\ref{sec:benchmark}.
We also audited a 357-record strict training split and ran clean phase-2
diagnostics; those diagnostics did not improve the main ScopeCtrl/NegRank
tradeoff, so the main results keep the strongest original adapter evaluated on
the audited test subset.

\subsection{Main Results}

Table~\ref{tab:main} reports results on the label-audited E4 v2 test subset
(397 records).

\begin{table*}[t]
\centering
\small
\caption{Main results on label-audited E4 v2 (397 records). NegRank is the
  single-gold Hard NegRank. BoolQ$^\dagger$: full validation
  (3,270 records); format \textbf{raw(PMI)} when null-prior No-bias $>0.5$.
  MMLU $\Delta$: lightweight 570-question diagnostic relative to base model.
  Best value in each column is \textbf{bold}; $\dagger$see \S\ref{sec:ood}.}
\label{tab:main}
\setlength{\tabcolsep}{5pt}
\begin{tabular}{lrrrrrrr}
\toprule
Method & NegRank & NegFlipAcc & ScopeCtrl & OverNeg & WikiFact & MMLU $\Delta$ & BoolQ$^\dagger$ \\
\midrule
\multicolumn{8}{l}{\textit{Qwen2.5-7B}} \\
Base (no fine-tuning) & 39.6 & 17.7 & 87.8 & 12.2 & --- & --- & 83.3 \\
\quad + Warning & 38.1 & \phantom{0}7.7 & 80.6 & 19.4 & --- & --- & --- \\
\quad + Persona & 37.6 & 12.0 & 77.6 & 22.4 & --- & --- & --- \\
\quad + CoT & 37.1 & 11.0 & 83.7 & 16.3 & --- & --- & --- \\
\quad + Contrastive Dec. & 36.6 & 20.7 & 71.4 & 28.6 & --- & --- & --- \\
Vanilla SFT & 59.9 & 61.2 & 54.1 & 45.9 & --- & $-0.88$ & --- \\
DPO & 44.6 & 25.8 & 76.5 & 23.5 & 44.0 & $+0.00$ & 86.8(87.0) \\
NC-SFT & 53.0 & 50.8 & 77.6 & 22.4 & 63.0 & $+0.18$ & 87.8 \\
MGNM (ours) & \textbf{63.9} & \textbf{64.2} & \textbf{92.9} & \phantom{0}\textbf{7.1} & \textbf{91.0} & $-1.23$ & \textbf{86.2(87.2)} \\
\midrule
\multicolumn{8}{l}{\textit{Llama-3.1-8B}} \\
Base & 30.7 & 12.0 & 81.6 & 18.4 & --- & --- & 83.2 \\
Vanilla SFT & 43.6 & 54.5 & 28.6 & 71.4 & --- & $-8.77$ & --- \\
DPO & 48.5 & 21.4 & 73.5 & 26.5 & 33.0 & $+1.23$ & 57.1(67.7) \\
NC-SFT & 58.4 & 53.2 & 68.4 & 31.6 & 75.0 & $-3.68$ & 85.5 \\
MGNM (ours) & \textbf{62.4} & \textbf{60.9} & \textbf{89.8} & \phantom{0}\textbf{10.2} & 79.0 & $-3.51$ & 69.0(81.4) \\
\midrule
\multicolumn{8}{l}{\textit{Mistral-7B (partial cross-architecture validation)}} \\
Base & 33.7 & 14.0 & 82.7 & 17.3 & --- & --- & --- \\
NC-SFT & 59.9 & 47.2 & 73.5 & 26.5 & 74.0 & $-0.53$ & --- \\
MGNM (ours) & 53.5 & 57.7 & 91.8 & \phantom{0}8.2 & 86.0 & $+0.70$ & --- \\
\bottomrule
\end{tabular}
\end{table*}

\paragraph{Inference-time methods.}
All prompt variants reduce ScopeCtrl below the base model (87.8\%$\to$78--84\%)
while yielding minimal NegFlipAcc gains (8--12\% vs.\ 17.7\% base).
CD achieves the highest inference-time NegFlipAcc (20.7\%) but collapses
ScopeCtrl to 71.4\% ($-16.4$pp), because amplifying the expert--amateur logit
gap increases differentiation without distinguishing suppress from preserve behavior.

\paragraph{Fine-tuning methods.}
Vanilla SFT achieves high NegFlipAcc (61.2\%) via wholesale suppression,
at the cost of severe over-negation (ScopeCtrl 54.1\%).
DPO and NC-SFT improve over base but reveal a fundamental tension:
suppression gains come at the cost of specificity.

\paragraph{MGNM.}
MGNM achieves 64.2\% NegFlipAcc and 92.9\% ScopeCtrl on Qwen2.5-7B,
with OverNeg at 7.1\%---\emph{lower} than the base model (12.2\%).
Improvements over DPO and NC-SFT remain statistically significant
(two-sided McNemar test): $+38.5$pp and $+13.4$pp on NegFlipAcc;
$+16.3$pp and $+15.3$pp on ScopeCtrl.
Results are consistent on Llama-3.1-8B, while Mistral-7B provides
partial cross-architecture validation.
On Mistral-7B-Instruct-v0.3, MGNM improves over its base model by
$+43.7$pp NegFlipAcc, $+9.1$pp ScopeCtrl, and $+19.8$pp NegRank,
while reducing OverNeg by $9.1$pp.
Compared with Mistral NC-SFT, MGNM is not highest on every single metric:
NC-SFT has higher hard NegRank (59.9 vs.\ 53.5). However, NC-SFT's ScopeCtrl
drops sharply to 73.5\% and OverNeg rises to 26.5\%, whereas MGNM preserves a
balanced behavior profile with stronger NegFlipAcc, ScopeCtrl, OverNeg, and
WikiFact.

\subsection{Open-Ended Generation Evaluation}
\label{sec:freegen}

Candidate-ranking evaluation may not fully reflect real-world usage where no
candidate pool is provided.
To complement the main results, we evaluate on the 129 \texttt{suppress\_target}
records using free-form greedy generation (80 tokens), then apply a
keyword-overlap heuristic ($\geq 60$\% overlap with banned entity keywords) as
an automatic proxy for whether the model reproduced the forbidden content
(Table~\ref{tab:freegen}). This result is supplementary and is not treated as
the primary benchmark.

\begin{table}[t]
\centering
\small
\caption{Supplementary free-generation evaluation on \texttt{suppress\_target} ($n{=}129$) and
  label-audited \texttt{preserve\_positive} ($n{=}98$) records.
  Suppress accuracy: keyword overlap with forbidden entity; Preserve accuracy:
  keyword overlap with gold-positive answer.}
\label{tab:freegen}
\setlength{\tabcolsep}{5pt}
\begin{tabular}{lrrrr}
\toprule
Model & \multicolumn{2}{c}{Suppress ($n{=}129$)} & \multicolumn{2}{c}{Preserve ($n{=}98$)} \\
\cmidrule(lr){2-3}\cmidrule(lr){4-5}
 & Free-gen & Cand-rank & Free-gen & Cand-rank \\
\midrule
Base       & 62.0\% & 17.7\% & 87.8\% & --- \\
DPO        & 62.8\% & 26.0\% & ---    & --- \\
NC-SFT     & 63.6\% & 52.0\% & ---    & --- \\
MGNM (oracle token)  & \textbf{96.1\%} & \textbf{64.2\%} & 85.7\% & \textbf{92.9\%} \\
MGNM (hidden ctrl.) & --- & --- & 85.7\% & --- \\
MGNM (no token)      & 69.0\% & --- & \textbf{92.9\%} & --- \\
\bottomrule
\end{tabular}
\end{table}

MGNM with oracle behavior token reaches 96.1\% free-gen NegFlipAcc on suppress
records under the keyword proxy, $+33$pp above the baselines (62--64\%).
This suggests that multi-grained training also improves open-ended generation,
but only as supplementary evidence; candidate ranking remains the main
benchmark.
Even without a behavior token at inference, MGNM retains a $+7$pp lift
(69.0\% vs.\ 62.0\% base), showing that training itself contributes limited
transfer beyond token-conditioned routing.
For audited preserve records, the base model achieves 87.8\% keyword overlap in free
generation---it often ignores the negating modifier and answers naturally---while
MGNM oracle reaches 85.7\%.
We further tested hidden-control variants that move the preserve instruction into
a system message, XML-style control tag, or natural-language instruction; the
best hidden-control setting remains 85.7\%, and plain-text rendering does not
improve this result.
Removing the \textsc{[Preserve]} prefix at inference raises preserve free-generation
to 92.9\%, indicating that the token is useful for candidate ranking but changes
the continuation distribution in open-ended generation in a way prompt-only
hidden control does not reliably repair.
Thus preserve free-generation is better decoded without a behavior prefix,
while the underlying candidate-ranking ScopeCtrl (92.9\%) remains strong.
The keyword-overlap detector is conservative for suppress; several of the failures
were cases where the model paraphrased the banned entity, so actual avoidance
rates may be higher.

\subsection{Out-of-Domain Generalization}
\label{sec:ood}

\paragraph{WikiFact.}
We evaluate on 200 negated factual assertions about Wikipedia entities
unseen during training.
MGNM achieves 91.0\% NegFlipAcc on Qwen (DPO: 44.0\%, NC-SFT: 63.0\%)
and 79.0\% on Llama (DPO: 33.0\%, NC-SFT: 75.0\%),
confirming that multi-grained training generalizes beyond template-based probes.

\paragraph{BoolQ.}
We evaluate on the BoolQ full validation set (3,270 yes/no questions~\cite{clark2019boolq}).
Qwen MGNM raw accuracy (86.2\%) is marginally \emph{higher} than base (83.3\%);
however, Llama MGNM (69.0\%) regresses substantially versus Llama base (83.2\%).
We investigate this Llama regression.

\emph{Prior-shift diagnosis.}
MGNM training shifts the model's null-context prior toward ``No.''
Formally, we measure $\Delta = \log P(\text{No}|\text{null}) - \log P(\text{Yes}|\text{null})$,
where the null prompt contains no passage.
Table~\ref{tab:prior} shows that MGNM and DPO induce
$\Delta \approx +1.0$ on both models, while NC-SFT and base show $|\Delta| < 0.15$.
This is a direct side-effect of $\mathcal{L}_\text{sup}$: reducing the log-probability
of affirmative completions globally shifts the null prior.

\begin{table}[t]
\centering
\small
\caption{Null-context prior shift $\Delta = \log P(\text{No}|\text{null}) - \log P(\text{Yes}|\text{null})$.
  Values $>0.5$ trigger PMI calibration.}
\label{tab:prior}
\begin{tabular}{lrr}
\toprule
Method & Qwen $\Delta$ & Llama $\Delta$ \\
\midrule
Base & $+0.09$ & $-0.68$ \\
DPO & $+1.09$ & $+1.09$ \\
NC-SFT & $+0.13$ & $-0.06$ \\
MGNM & $+1.06$ & $+0.80$ \\
\bottomrule
\end{tabular}
\end{table}

\emph{PMI calibration.}
We apply PMI~\cite{holtzman2021surface}:
$\hat{s}(p, c) = s(p, c) - \log P(c \mid \text{null})$.
After calibration, Qwen MGNM achieves \textbf{87.2\%}, the highest PMI-calibrated
accuracy across all evaluated methods.
Llama MGNM reaches 81.4\%, a residual 1.8pp gap versus Llama base (83.2\% raw).
The large Llama raw regression is therefore best interpreted as a combination
of calibration shift and training tradeoff, rather than as a broad capability
collapse. A stricter held-out calibration diagnostic supports the calibration
component: using a 20\% calibration split
and 80\% held-out test split over five random seeds, the best scalar PMI
coefficient is stable at $\alpha = 1.65 \pm 0.12$, raising held-out Llama MGNM
accuracy from 81.4\% under standard PMI to 84.6\% on average. At the same time,
retention-oriented checkpoints can recover raw BoolQ only by weakening E4
ranking, ScopeCtrl, or WikiFact, so we report BoolQ as a calibration-plus-training
tradeoff rather than calibration alone. Qwen shows no regression
even in raw accuracy.

\subsection{Ablation Study}

Table~\ref{tab:ablation} isolates each loss component.

\begin{table}[t]
\centering
\small
\caption{Ablation study (E4 v2). $\Delta$ relative to MGNM full.}
\label{tab:ablation}
\setlength{\tabcolsep}{4pt}
\begin{tabular}{lrrrr}
\toprule
Variant & NegFlipAcc & ScopeCtrl & OverNeg & NegRank \\
\midrule
\multicolumn{5}{l}{\textit{Qwen2.5-7B}} \\
MGNM (full) & \textbf{64.2} & \textbf{92.9} & \phantom{0}\textbf{7.1} & \textbf{63.9} \\
w/o $\mathcal{L}_\text{sup}$ & 31.8 & 86.7 & 13.3 & 48.5 \\
w/o $\mathcal{L}_\text{pre}$ & 49.2 & 13.3 & 86.7 & 36.1 \\
\midrule
\multicolumn{5}{l}{\textit{Llama-3.1-8B}} \\
MGNM (full) & \textbf{60.9} & \textbf{89.8} & \phantom{0}\textbf{10.2} & \textbf{62.4} \\
w/o $\mathcal{L}_\text{sup}$ & 44.8 & 81.6 & 18.4 & 48.5 \\
w/o $\mathcal{L}_\text{pre}$ & 55.2 & 22.4 & 77.6 & 44.1 \\
\bottomrule
\end{tabular}
\end{table}

Removing $\mathcal{L}_\text{sup}$ drops NegFlipAcc by $-32.4$pp (Qwen)
and $-16.1$pp (Llama), confirming it is the primary driver of
suppression accuracy.
Removing $\mathcal{L}_\text{pre}$ is catastrophic for ScopeCtrl:
$-79.6$pp (Qwen) and $-67.3$pp (Llama), with OverNeg rising sharply---
mirroring Vanilla SFT.
Both components are individually necessary: $\mathcal{L}_\text{sup}$ trains
\emph{what} to suppress; $\mathcal{L}_\text{pre}$ trains \emph{when not to}.

\paragraph{Behavior token ablation.}
To quantify oracle-token dependence, we evaluate the same MGNM model under
multiple token conditions, including structure-assisted routing strategies, and
additionally train a no-token variant (Table~\ref{tab:token_ablation}).

\begin{table}[t]
\centering
\small
\caption{Behavior token ablation (Qwen2.5-7B, candidate ranking, $n{=}397$).
  Router accuracy denotes the fraction of records assigned the correct token.}
\label{tab:token_ablation}
\setlength{\tabcolsep}{4pt}
\begin{tabular}{lrrrl}
\toprule
Setting & NegFlipAcc & ScopeCtrl & NegRank & Note \\
\midrule
Oracle token (current)         & \textbf{64.2} & \textbf{92.9} & \textbf{63.9} & correct token at inference \\
\midrule
Two-stage + suppression-only guard & 62.9 & 88.8 & 61.4 & force \texttt{suppression\_only} to suppress \\
Two-stage supervised router (92.9\% acc) & 61.9 & 87.8 & 60.4 & preserve-first classifier \\
Two-stage supervised, empty SELECT & 61.9 & 88.8 & 61.4 & no prefix when classifier predicts select \\
Candidate-aware LLM router (69.5\% acc) & 60.5 & 66.3 & 59.9 & prompt plus candidate pool \\
Candidate LLM, empty SELECT    & 59.5 & 69.4 & 61.4 & no prefix when router predicts select \\
Supervised TF-IDF router (70.0\% acc) & 56.9 & 53.1 & 57.4 & trained on strict-train labels \\
Supervised TF-IDF, empty SELECT & 57.2 & 59.2 & 59.4 & no prefix when classifier predicts select \\
Prompt-only LLM router         & 42.1 & 73.5 & 40.6 & GPT-4.1-mini zero-shot routing \\
Rule router (67.3\% acc)       & 53.5 & 44.9 & 45.5 & keyword/regex heuristics \\
\midrule
No token (oracle model)        & 59.5 & 56.1 & \textbf{64.9} & no prefix at inference \\
Wrong token (all \textsc{[Suppress]}) & 42.8 & \phantom{0}0.0 & 24.8 & adversarial mismatch \\
No-token trained model         & 57.2 & 67.3 & 62.4 & trained without prefixes \\
\bottomrule
\end{tabular}
\end{table}

Without any token at inference, NegFlipAcc drops $-4.7$pp while ScopeCtrl falls $-36.7$pp;
NegRank is slightly higher ($+1.0$pp) because \texttt{select\_gold\_neg}
records already have no token in training.
Under the adversarial wrong-token condition (forcing \textsc{[Suppress]} on all
records), ScopeCtrl collapses to $0.0\%$, confirming tokens carry genuine
behavioral semantics rather than serving as inert prefixes.
Training without any behavior prefix yields NegFlipAcc $57.2\%$ ($-7.0$pp) and
ScopeCtrl $67.3\%$ ($-25.6$pp), showing that behavior tokens contribute at
\emph{both} training and inference: without them at training time, the model
cannot fully learn scope-conditional behavior, raising over-negation to $32.7\%$
on preserve records despite retaining partial suppression capability.

\paragraph{Structure-assisted routing.}
We explore structure-assisted token-assignment strategies ranging from keyword/regex
rules to a candidate-aware LLM router (GPT-4.1-mini).
Adding the unlabeled candidate pool raises router accuracy to 69.5\% and
improves downstream performance to 60.5\% NegFlipAcc, 66.3\% ScopeCtrl, and
59.9\% NegRank.
Because \texttt{select\_gold\_neg} rows are relatively robust without a prefix,
using an empty prefix when the router predicts SELECT further improves the
Scope/Rank tradeoff (59.5/69.4/61.4).
Error analysis shows that the main failure mode is preserve cases being
over-routed to SELECT. We therefore train a two-stage supervised router:
stage 1 detects \texttt{preserve\_positive} versus non-preserve, and stage 2
separates \texttt{suppress\_target} from \texttt{select\_gold\_neg}. With a
threshold selected on a held-out calibration split, this router reaches 92.9\%
behavior accuracy on the strict test subset (suppress 94.8\%, preserve 93.9\%,
select 91.6\%). Downstream, it achieves 61.9\% NegFlipAcc, 87.8\% ScopeCtrl,
and 60.4\% NegRank; using an empty prefix when the router predicts SELECT
further improves the automatic setting to 61.9/88.8/61.4. Adding a
\texttt{suppression\_only} guard and keeping the calibrated threshold at 0.35
raises behavior accuracy to 94.2\% and downstream performance to
62.9/88.8/61.4. We also test a text-derived suppression guard: a small set of
surface-pattern detectors routes explicit exclusion questions to \suppress{}
without reading dataset-internal \texttt{semantic\_mode} labels. Table~\ref{tab:router_guard}
shows that this text-derived guard recovers the WikiFact suppress-only failure
to the oracle-token level.
These results define the scope of our claim: MGNM is a token-conditioned
controller whose peak performance assumes correct behavior selection, while
automatic routing remains a deployment bottleneck.

\begin{table}[t]
\centering
\small
\caption{Suppression guard variants for automatic routing. E4 scores are
  NegFlipAcc / ScopeCtrl / NegRank on the 397-record strict test subset.
  WikiFact is a 100-record suppress-only out-of-domain probe.}
\label{tab:router_guard}
\setlength{\tabcolsep}{3pt}
\begin{tabular}{lrr}
\toprule
Router & E4 & WikiFact recall / NegFlipAcc \\
\midrule
Two-stage router & 61.9 / 87.8 / 60.4 & 29.0 / 78.0 \\
+ metadata guard & 62.9 / 88.8 / 61.4 & 100.0 / 91.0 \\
+ text-derived guard & 61.9 / 88.8 / 61.4 & 100.0 / 91.0 \\
\bottomrule
\end{tabular}
\end{table}

\paragraph{SELECT multi-answer audit.}
Since some SELECT prompts admit multiple semantically valid negative answers,
we report both single-gold Hard NegRank and audited Multi-answer NegRank.
Hard NegRank is intentionally strict and accepts only the original gold
negative answer. Multi-answer NegRank additionally accepts manually verified
valid negative alternatives. For example,
when a question asks which country does not primarily speak Spanish, both
``Germany'' and ``Japan'' may be valid negative candidates depending on the
candidate pool. We therefore run a conservative manual audit over candidate-pool
misses and report Multi-answer NegRank in Table~\ref{tab:select_audit}.

\begin{table}[t]
\centering
\small
\caption{SELECT single-gold ambiguity audit on \texttt{select\_gold\_neg}
records ($n{=}202$). Multi-answer NegRank counts manually verified valid
alternatives from \texttt{candidate\_pool\_neg}.}
\label{tab:select_audit}
\setlength{\tabcolsep}{5pt}
\begin{tabular}{lrrr}
\toprule
Model & Hard NegRank & Multi-answer NegRank & Rescued \\
\midrule
Qwen Base & 36.6 & 45.0 & 17 \\
Qwen DPO & 44.6 & 55.9 & 23 \\
Qwen NC-SFT & 53.0 & 64.4 & 23 \\
Qwen MGNM & \textbf{63.9} & \textbf{85.6} & 44 \\
\bottomrule
\end{tabular}
\end{table}

For Qwen MGNM, 44 of 50 candidate-pool misses are judged valid under this
conservative audit, raising SELECT NegRank from 63.9\% to 85.6\%. This means
many hard misses are annotation-boundary cases rather than invalid model
choices. Applying the same audit protocol to all Qwen models preserves MGNM's
advantage and indicates that the main remaining SELECT bottleneck is partly a
label-boundary issue, not simply insufficient rank-loss strength. As a check, a stronger rank-dist
variant (RD30) does not improve full-test NegRank and worsens ScopeCtrl, so we
keep the current MGNM model as the main result.

\subsection{Data Scale Analysis}

Figure~\ref{fig:scale} shows NegFlipAcc, ScopeCtrl, and Hard NegRank as a function
of training set fraction, with random 25\%/50\%/75\%/100\% samples from
the 471-record augmented training input (same seed for all, oversampling preserved).

\begin{figure}[t]
\centering
\includegraphics[width=0.95\columnwidth]{../outputs/data_scale_curve.pdf}
\caption{Performance vs.\ training data scale (Qwen2.5-7B).
  Marginal gains decrease after 50\%; 25\% of training data already outperforms
  all inference-time methods, showing high sample efficiency.}
\label{fig:scale}
\end{figure}

The learning curve is notably steep at low data: after reaggregation on the
audited test subset, 50\% of training data (236 records) yields 53.8\% NegFlipAcc
and 85.7\% ScopeCtrl.
Gains from 50\%$\to$100\% are smaller than the early-data gains
($+10.4$pp NegFlipAcc, $+7.1$pp ScopeCtrl),
and 25\% of training data already exceeds all inference-time methods on NegFlipAcc,
indicating high sample efficiency.
Whether scaling the dataset further continues to improve performance on more
diverse distributions remains an open question.

%% ─────────────────────────────────────────────────────────────────────────────
\section{Mechanistic Analysis}
\label{sec:mechanism}

To understand \emph{where} negation blindness originates, we apply three
mechanistic probes to 200 E4 probes: logit-lens (E1),
attention attribution (E2), and activation patching (E3).

\begin{figure}[t]
\centering
\includegraphics[width=0.95\columnwidth]{../outputs/mechanism_cross_model_final.pdf}
\caption{Mechanistic comparison: Qwen2.5-7B Base, Qwen MGNM, and Llama MGNM.
  \textbf{Top}: neg\_margin = allowed $-$ forbidden score (logit lens).
  \textbf{Middle left}: best-allowed rate; \textbf{Middle right}: attention to negation token.
  \textbf{Bottom}: activation patching recovery rate.}
\label{fig:mechanism}
\end{figure}

\paragraph{E1: Layer-wise negation signal.}
We compute the \emph{neg\_margin} = $\max_{c^+ \in \mathcal{G}^+}
\hat{s}_\ell(p^-, c^+) - s_\ell(p^-, c^*)$ at each layer $\ell$
via the logit lens~\cite{nostalgebraist2020logitlens},
where $c^*$ is the forbidden target.
A negative margin indicates the forbidden target outranks allowed alternatives.
As shown in Figure~\ref{fig:mechanism} (top), the \emph{mean} margin (across 200 probes)
is negative throughout all 28 layers in the Qwen base model (last-layer mean: $-2.52$).
MGNM modestly improves the final-layer best-allowed rate from 0.535 to 0.595.
Note: E1 reports layer-wise means, while NegFlipAcc is a sample-level strict-ranking metric;
the base model's 17.7\% NegFlipAcc indicates that $\sim$18\% of audited examples individually show
allowed candidates outranking the forbidden target, consistent with a negative mean margin.

\paragraph{E2: Attention to negation tokens.}
The final-position token allocates substantial attention to negation tokens
(max-head weight $\approx0.80$ at last layer), consistent across base and MGNM.
This rules out \emph{attention blindness} as the cause: the model attends to
negation tokens but does not propagate the negation signal.

\paragraph{E3: Activation patching.}
We causally patch activations from an unnegated (positive) counterpart prompt
(\emph{clean}) into the negated prompt's processing (\emph{corrupt}), replacing
layer-$\ell$ activations and measuring how much correct-ranking recovery results.
Peak positive-recovery rate is 55\% (base) and 60\% (MGNM) when patching
from \emph{layer 0}, declining monotonically with patch layer.
This suggests the clean representation is most informative at layer 0 and is
progressively overwritten by later context; however, this interpretation
depends on the patching design and should be treated as suggestive rather than
conclusive evidence of early-layer overwriting.

\paragraph{Interpretation.}
The large NegFlipAcc gain on the audited subset ($+46.5$pp) is disproportionate to the small
representational shift ($+6$pp best-allowed rate at last layer).
This suggests MGNM operates through two complementary mechanisms:
(1) partial repair of internal negation representations, and
(2) behavior-token conditioning that re-routes output decisions at
the final layer, independent of intermediate representations.
This interpretation explains why inference-time interventions fail:
they can only modify (2), but not (1).

%% ─────────────────────────────────────────────────────────────────────────────
\section{Analysis}
\label{sec:analysis}

\subsection{Why Inference-Time Methods Fail}

Our mechanistic results explain the failure pattern:
because the negation signal often fails to become a candidate-level ranking
preference (with E3's largest recovery at layer 0), downstream logit
manipulation is poorly positioned to recover it.
CD amplifies the expert--amateur logit difference, which reflects semantic
plausibility rankings, not negation scope, causing suppress and preserve decisions
to be handled identically.
CoT reasoning ($-7.2$pp NegFlipAcc vs.\ base) and persona prompting ($-5.7$pp ScopeCtrl)
introduce surface-level behavior that does not alter the probability assignments
governing final token selection.
MGNM vs.\ CD on the audited subset: NegFlipAcc $+43.5$pp and ScopeCtrl $+21.5$pp.

\subsection{Prior Shift as a Diagnostic}

The null-context prior shift $\Delta$ is a strong predictor of BoolQ raw degradation
\emph{when the base model has a large natural bias}.
For Llama (base $\Delta = -0.68$, Yes-biased), methods with $\Delta > 0.5$ (MGNM: $+0.80$,
DPO: $+1.09$) flip the prior to No-biased, causing a large raw accuracy drop.
For Qwen (base $\Delta = +0.09$, near-neutral), the same $\Delta > 0.5$ shift
produces mild raw \emph{improvements} (86.2\% and 86.8\% vs.\ 83.3\% base)
because the prior shift corrects a pre-existing slight Yes-bias.
$\Delta$ thus identifies models at risk of BoolQ degradation: a pre-training
Yes-bias ($\Delta_\text{base} < 0$) combined with a post-training $\Delta > 0.5$
warrants PMI calibration.
Models without PMI calibration applied are expected to degrade.

\subsection{Qualitative Case Study}

Table~\ref{tab:cases} shows which candidate each model ranks first for two
prompts, illustrating the suppress--preserve asymmetry that prior methods collapse.

\begin{table}[t]
\centering
\small
\caption{Top-ranked candidate per model for two E4 records.
  (\textbf{bold} = correct; \textit{italic} = failure; \checkmark/\texttimes{} = pass/fail).}
\label{tab:cases}
\setlength{\tabcolsep}{3pt}
\begin{tabular}{lll}
\toprule
Model & Top candidate & Pass? \\
\midrule
\multicolumn{3}{l}{\textit{Case A (suppress): ``Is metal NOT known for high conductivity?''}} \\
\multicolumn{3}{l}{\quad Banned: \textit{``Metal is known for its high conductivity.''}} \\
Base   & \textit{Metal is known for its high conductivity.} & \texttimes \\
DPO    & \textit{Metal is known for its high conductivity.} & \texttimes \\
NC-SFT & \textit{Metal is known for its high conductivity.} & \texttimes \\
MGNM   & \textbf{Metal is known for its high strength.}     & \checkmark \\
\midrule
\multicolumn{3}{l}{\textit{Case B (preserve): ``Do not skip preheating the oven before baking the cake. Why is this important?''}} \\
\multicolumn{3}{l}{\quad Preserved answer: ``Preheating ensures the cake bakes evenly and rises properly.''} \\
Base   & \textit{Use a toothpick to test if the cake is done.} & \texttimes \\
DPO    & \textit{Use a toothpick to test if the cake is done.} & \texttimes \\
NC-SFT & \textit{Use a toothpick to test if the cake is done.} & \texttimes \\
MGNM   & \textbf{Preheating ensures the cake bakes evenly and rises properly.} & \checkmark \\
\bottomrule
\end{tabular}
\end{table}

In Case A, all baselines output the forbidden answer despite the explicit
\emph{NOT} constraint; MGNM with \suppress{} token correctly suppresses it.
In Case B, the negation scopes over \emph{skipping} preheating, so the factual
answer about why preheating matters is still entailed and should be
preserved; all baselines are misled by the negated surface form and rank an
irrelevant candidate first, while MGNM with \preserve{} token retains the
correct affirmative answer.

%% ─────────────────────────────────────────────────────────────────────────────
\section{Limitations}
\label{sec:limitations}

\paragraph{Inference-time behavior selection.}
MGNM's strongest E4 results assume that the caller supplies the correct
behavior token. Removing the token preserves some suppression ability
(59.5\% NegFlipAcc) but reduces ScopeCtrl from 92.9\% to 56.1\%.
A candidate-aware LLM router improves the routed setting to 60.5\% NegFlipAcc,
66.3\% ScopeCtrl, and 59.9\% NegRank, and an empty-prefix SELECT variant
improves the Scope/Rank tradeoff to 59.5/69.4/61.4.
A two-stage supervised router that first detects preserve records reaches
92.9\% behavior classification accuracy and improves the automatic setting to
61.9\% NegFlipAcc, 88.8\% ScopeCtrl, and 61.4\% NegRank when SELECT predictions
use an empty prefix. Adding a \texttt{suppression\_only} guard raises strict-test
behavior accuracy to 94.2\% and downstream performance to 62.9/88.8/61.4.
On WikiFact, a suppress-only out-of-domain probe, the unguarded router recovers
only 29.0\% suppress recall and reduces Qwen MGNM from 91.0\% oracle-token
NegFlipAcc to 78.0\%. Both the metadata guard and a text-derived suppression
detector restore 100.0\% suppress recall and 91.0\% downstream NegFlipAcc.
This shows that explicit exclusion prompts can be handled without relying on
dataset-internal \texttt{semantic\_mode} labels. However, arbitrary pure-text
behavior inference remains open. MGNM should therefore be interpreted as a
conditional negation behavior controller, not as a complete domain-general
automatic router.

\paragraph{Behavior token task specificity.}
The \preserve{} token is effective within the E4 task format but does not
transfer to other tasks.
Adding \preserve{} to BoolQ prompts causes Qwen MGNM accuracy to drop to 47\%
(near-random), showing that behavior tokens are coupled to the training input
distribution.
Extending behavior token conditioning to arbitrary tasks would require
multi-task training or a prompt-adaptation step.

\paragraph{Residual Llama BoolQ gap.}
After PMI calibration, Llama MGNM (81.4\%) trails Llama base raw (83.2\%)
by 1.8pp under the default $\alpha=1$ protocol. However, a held-out calibration
diagnostic shows that the calibration component is recoverable: selecting a scalar PMI
coefficient on a separate 20\% calibration split yields a stable optimum around
$\alpha=1.65$ and 84.6\% average held-out accuracy across five seeds. We
therefore treat calibration shift as one important factor. The retention
diagnostics show the other factor: checkpoints that recover raw BoolQ do so by
weakening E4 ranking, ScopeCtrl, or WikiFact. We therefore interpret Llama as
revealing a real training tradeoff between yes/no retention and specialized
negation-control behavior, not as a simple model failure. For a formal
paper-facing result, any scalar calibration still requires a single pre-declared
calibration split rather than repeated split selection.

\paragraph{Scope of negation phenomena.}
E4 is designed as a controlled diagnostic benchmark rather than a large-scale
coverage benchmark: its value is isolating behavior types and enabling causal
comparison under strict holdout. It covers lexical and morphological negation
in entity-selection contexts.
Compositional negation, presupposition cancellation, and multi-hop negation
are outside the current benchmark scope. The audited strict training split has
only 357 records, while the strongest Qwen adapter uses the 471-record augmented
training input; larger audited training sets are needed to test whether clean-from-scratch
training can match or exceed the current strongest adapter.

%% ─────────────────────────────────────────────────────────────────────────────
\section{Conclusion}
\label{sec:conclusion}

We have presented MGNM, a LoRA fine-tuning approach that addresses negation
blindness in LLMs through caller-specified behavior tokens and five
complementary training objectives.
On the label-audited E4 benchmark, MGNM simultaneously achieves high suppression
accuracy (NegFlipAcc 64.2\%), high specificity (ScopeCtrl 92.9\%), and
over-negation suppression (OverNeg 7.1\%)---objectives that previous methods trade off
against each other.
Mechanistic probes (200 E4 candidate-ranking instances) suggest that negation
failure under this setting is not primarily an attention deficiency, but a failure
to translate the negation signal into candidate-level ranking preferences.
MGNM partially repairs this internal signal while additionally exposing
a conditional interface for output behavior, explaining its disproportionate
behavioral improvement relative to the representational shift.
The remaining deployment question is automatic behavior selection under known
task structure, which our router experiments show is substantially harder than
the conditional control problem solved here.
We will release E4, all training scripts, and model adapters upon publication.
For anonymous review, we provide anonymized data samples, evaluation scripts,
and metric-computation code.

%% ─────────────────────────────────────────────────────────────────────────────
\bibliography{references}

%% ─────────────────────────────────────────────────────────────────────────────
\appendix

\section{Reproducibility Statement}

All experiments load locally prepared copies of publicly available base models:
Qwen2.5-7B~\cite{qwen2025qwen25},
Llama-3.1-8B~\cite{dubey2024llama},
and Mistral-7B-Instruct-v0.3~\cite{jiang2023mistral}.
LoRA fine-tuning requires a single 80GB A100 GPU; training time is 15--40 minutes
(3 epochs on the 471-record augmented training input).
E4 evaluation on 397 audited records requires $<5$ minutes per model.
All code, data, and trained adapters will be released upon publication under
CC-BY-4.0 (dataset) and Apache-2.0 (code/adapters); anonymized data samples,
evaluation scripts, and metric computation code are provided for review.

Hyperparameter details are provided in \S\ref{sec:method}.
The complete training script exposes all hyperparameters as CLI arguments;
the exact invocation for each reported model is included in the supplementary
material's \texttt{run\_*.sh} scripts.

\paragraph{MMLU evaluation protocol.}
We use the standard Hendrycks et al.~\cite{hendrycks2020measuring} MMLU benchmark
(57 subjects). We sample 10 questions per subject from the validation split
(570 questions total, seed~42) and score each option via 5-shot log-probability
of the single-token continuation (A/B/C/D) following \texttt{Answer:}.
We use a custom scorer (\texttt{scripts/evaluate\_mmlu.py}), not lm-eval-harness.
The Llama MGNM $-3.51$pp drop corresponds to $\approx$20 questions out of 570
and reflects Llama-specific sensitivity to LoRA fine-tuning.
Because we evaluate a 570-question stratified subset rather than the full
MMLU test set, MMLU $\Delta$ should be interpreted as a lightweight
capability-retention diagnostic rather than a definitive estimate of
general knowledge degradation.

\section{Ethics Statement}

The E4 dataset is synthetically generated from Wikipedia entity facts and
procedural templates; no private data or crowdsourced human annotations are used.
Negation blindness mitigation is a safety-relevant capability: models that
ignore negation may produce outputs that violate user constraints in
high-stakes applications.
We identify no dual-use risks specific to negation handling.
Total compute for all reported experiments is approximately 4 GPU-hours on
an A100-80GB, excluding baseline model downloads.

\section{Full Significance Test Results}

\begin{table}[h]
\centering
\small
\caption{McNemar two-sided exact test, label-audited E4 v2.
  $b$~=~baseline-only-correct; $c$~=~MGNM-only-correct.
  Metric subsets: NegFlipAcc ($n{=}299$, \texttt{flip\_required}), ScopeCtrl ($n{=}98$, preserve), NegRank ($n{=}202$, select).}
\label{tab:sig}
\begin{tabular}{llrrrrrl}
\toprule
Comparison & Metric & MGNM & Base & $\Delta$ & $b$ & $c$ & Sig. \\
\midrule
Qwen vs.\ DPO & NegFlipAcc & 64.2 & 25.8 & $+38.5$ & 6 & 121 & *** \\
 & ScopeCtrl & 92.9 & 76.5 & $+16.3$ & 2 & 18 & *** \\
 & NegRank & 63.9 & 44.6 & $+19.3$ & 8 & 47 & *** \\
\midrule
Qwen vs.\ NC-SFT & NegFlipAcc & 64.2 & 50.8 & $+13.4$ & 16 & 56 & *** \\
 & ScopeCtrl & 92.9 & 77.6 & $+15.3$ & 3 & 18 & ** \\
 & NegRank & 63.9 & 53.0 & $+10.9$ & 12 & 34 & *** \\
\midrule
Llama vs.\ DPO & NegFlipAcc & 60.9 & 21.4 & $+39.5$ & 9 & 127 & *** \\
 & ScopeCtrl & 89.8 & 73.5 & $+16.3$ & 3 & 19 & *** \\
 & NegRank & 62.4 & 48.5 & $+13.9$ & 11 & 39 & *** \\
\midrule
Llama vs.\ NC-SFT & NegFlipAcc & 60.9 & 53.2 & $+7.7$ & 33 & 56 & * \\
 & ScopeCtrl & 89.8 & 68.4 & $+21.4$ & 6 & 27 & *** \\
 & NegRank & 62.4 & 58.4 & $+4.0$ & 21 & 29 & ns \\
\bottomrule
\multicolumn{8}{l}{*$p{<}0.05$, **$p{<}0.01$, ***$p{<}0.001$; ns~$=$~not significant ($p{=}0.16$, two-sided).} \\
\end{tabular}
\end{table}

\section{Prompt Templates}

The four prompt-based baselines use the following system messages
(truncated for space; full templates in supplementary):

\begin{description}
\item[Warning:] \emph{``Important: Strictly respect all negations and
  exclusions in the user's request. If the user says `not X', never output X.''}
\item[Persona:] \emph{``You are a precise language model with perfect negation
  comprehension. You always interpret `not', `never', `without' as strict
  exclusions.''}
\item[CoT:] \emph{``Step 1: Identify negation words. Step 2: List excluded concepts.
  Step 3: Generate candidates. Step 4: Remove any excluded concepts.
  Step 5: Output the final answer.''}
\end{description}

\end{document}
